% Plantilla inspirada en el cartel Seminario Chichimeca
% Compilar en VS Code con pdfLaTeX (LaTeX Workshop)
\documentclass[final]{beamer}
\usepackage[size=custom,width=84.1,height=118.9,scale=1.0]{beamerposter}
\usepackage[utf8]{inputenc}
\usepackage[T1]{fontenc}
\usepackage[spanish]{babel}
\usepackage{amsmath,amssymb,amsthm,mathtools}
\usepackage{tikz}
\usepackage{graphicx}
\usepackage{ragged2e}
\usepackage{lmodern}

% ===== COLORES EDITABLES =====
\definecolor{grisFondo}{HTML}{FFFFFF}
\definecolor{grisCabecera}{HTML}{EEEEEE}
\definecolor{grisSeccion}{HTML}{EEEEEE}
\definecolor{azulTitulo}{HTML}{27394C}
\definecolor{rojoAcento}{HTML}{B64A4B}
\definecolor{doradoAcento}{HTML}{A47C3C}
\definecolor{texto}{HTML}{333333}
\setbeamercolor{background canvas}{bg=grisFondo}
\setbeamercolor{normal text}{fg=texto,bg=grisFondo}
\setbeamertemplate{navigation symbols}{}
\setbeamertemplate{itemize item}{\color{azulTitulo}\raisebox{0.2ex}{\tiny$\blacksquare$}}
\setlength{\parindent}{0pt}
\setlength{\parskip}{0.5em}

% ===== DATOS DEL CARTEL =====
\newcommand{\TituloCartel}{Propiedades de separación\\inducidas por cocientes compactos}
\newcommand{\AutorCartel}{Juan Carlos Monter Cortés}
\newcommand{\InstitucionCartel}{Universidad de Guadalajara}
\newcommand{\EventoCartel}{Seminario Chichimeca de Álgebra}
\newcommand{\CorreoCartel}{juan.monter2902@alumnos.udg.mx}
\newcommand{\ProgramaCartel}{Doctorado en Ciencias Matemáticas}

% Logos opcionales: coloca logo_izq.png y logo_der.png en la carpeta
\newcommand{\LogoOpcional}[1]{%
  \IfFileExists{#1}{\includegraphics[width=8cm,height=8cm,keepaspectratio]{#1}}{%
  \begin{tikzpicture}\draw[rounded corners=3pt,draw=gray!55,dashed] (0,0) rectangle (8,8);
  \node[align=center,text=gray,font=\fontsize{16}{18}\selectfont] at (4,4) {LOGOTIPO\\institucional};\end{tikzpicture}}}

% ===== ENCABEZADO DE SECCION =====
\newcommand{\Seccion}[1]{%
  \vspace{0.6cm}
  \begin{beamercolorbox}[wd=\linewidth,sep=0.35cm,center]{section heading}
    {\color{azulTitulo}\bfseries\fontsize{22}{26}\selectfont #1}
  \end{beamercolorbox}
  \vspace{-0.2cm}\color{gray!75}\rule{\linewidth}{1.1pt}\color{texto}\vspace{0.25cm}
}
\setbeamercolor{section heading}{bg=grisSeccion,fg=azulTitulo}
\newcommand{\Subtitulo}[1]{\vspace{0.2cm}{\bfseries\color{azulTitulo} #1}\par}
\newcommand{\Destacado}[1]{{\color{rojoAcento}\bfseries #1}}
\newcommand{\Caja}[1]{%
 \begin{tikzpicture}\node[draw=doradoAcento,rounded corners=4pt,inner sep=10pt,text width=0.91\linewidth,align=left] {#1};\end{tikzpicture}}

\begin{document}
\begin{frame}[t]
\begin{tikzpicture}[remember picture,overlay]
  \fill[grisCabecera] (current page.north west) rectangle ([yshift=-12.7cm]current page.north east);
  \draw[gray!45,line width=0.6pt] ([yshift=-12.7cm]current page.north west) -- ([yshift=-12.7cm]current page.north east);
\end{tikzpicture}
\vspace{-0.4cm}
\begin{columns}[c,totalwidth=0.96\paperwidth]
 \begin{column}{0.13\paperwidth}\centering\LogoOpcional{logo_izq.png}\end{column}
 \begin{column}{0.70\paperwidth}\centering
  {\color{black}\bfseries\fontsize{36}{42}\selectfont \TituloCartel\par}
  \vspace{0.45cm}{\fontsize{23}{28}\selectfont \AutorCartel\par}
  {\fontsize{19}{23}\selectfont \InstitucionCartel\par}
 \end{column}
 \begin{column}{0.13\paperwidth}\centering\LogoOpcional{logo_der.png}\end{column}
\end{columns}
\vspace{2.3cm}

\begin{columns}[t,totalwidth=0.94\paperwidth]
\begin{column}{0.455\paperwidth}
\fontsize{19}{24}\selectfont\justifying
\Seccion{Introducción y motivación}
Aquí puedes presentar el problema central de tu investigación. La plantilla conserva el estilo de dos columnas y encabezados grises del cartel de referencia.

\Seccion{Preliminares}
\Subtitulo{Definición}
Un \emph{marco} es un retículo completo $A$ en el que se cumple
\[
 a\wedge\bigvee S=\bigvee\{a\wedge s:s\in S\}.
\]
\Subtitulo{Notación}
Escribe aquí tus conceptos básicos, definiciones y referencias.

\Seccion{Resultados principales}
\Subtitulo{Teorema A (enunciado de ejemplo)}
\Caja{Sea $A$ un marco que satisface las hipótesis indicadas. Entonces se obtiene la siguiente propiedad de separación.}

\Destacado{Idea de la demostración.} Añade aquí el argumento, un esquema o una ecuación relevante.
\[
  A\xrightarrow{\ f\ }B,\qquad f_*:B\longrightarrow A.
\]

\Seccion{Ejemplos y observaciones}
Puedes incluir contraejemplos, construcciones particulares o diagramas de implicaciones.
\end{column}
\hspace{0.02\paperwidth}
\begin{column}{0.455\paperwidth}
\fontsize{19}{24}\selectfont\justifying
\Seccion{Cocientes compactos}
\Subtitulo{Planteamiento}
Explica aquí el papel de los cocientes compactos y su relación con los axiomas de separación en la categoría de marcos.

\Seccion{Construcción o metodología}
\Subtitulo{Paso 1.} Presenta la construcción.
\Subtitulo{Paso 2.} Describe sus propiedades.
\Subtitulo{Paso 3.} Expón la conclusión.
\[
\begin{tikzpicture}[>=stealth,baseline]
\node (A) at (0,0) {$A$};\node (B) at (5,0) {$A_j$};
\draw[->,thick,azulTitulo] (A)--node[above] {$q$}(B);
\end{tikzpicture}
\]

\Seccion{Consecuencias}
\Subtitulo{Corolario (ejemplo)}
Escribe aquí una consecuencia del teorema principal. Puedes resaltar una afirmación breve con el color de acento.

\Destacado{Observación.} Los títulos y las dimensiones se modifican en la sección de configuración del archivo.

\Seccion{Conclusiones}
Resume las aportaciones principales y menciona las preguntas abiertas.

\Seccion{Referencias}
{\fontsize{16}{20}\selectfont
[1] P. T. Johnstone, \emph{Stone Spaces}, Cambridge University Press, 1982.\par
[2] J. Picado y A. Pultr, \emph{Frames and Locales}, Springer, 2012.\par
[3] J. Picado y A. Pultr, \emph{Separation in Point-Free Topology}, Springer, 2021.\par}
\end{column}
\end{columns}
\vfill
\color{gray!60}\rule{0.94\paperwidth}{0.6pt}
\begin{columns}[c,totalwidth=0.94\paperwidth]
\begin{column}{0.32\paperwidth}{\fontsize{15}{19}\selectfont \ProgramaCartel}\end{column}
\begin{column}{0.36\paperwidth}\centering{\fontsize{15}{19}\selectfont \EventoCartel}\end{column}
\begin{column}{0.32\paperwidth}\raggedleft{\fontsize{15}{19}\selectfont \CorreoCartel}\end{column}
\end{columns}
\vspace{0.5cm}
\end{frame}
\end{document}
